\documentclass[10pt,a4paper]{amsart}
\usepackage[margin=19mm]{geometry}
\usepackage{amsmath,amssymb,mathtools}
\usepackage[scr=rsfs,cal=euler]{mathalfa}
\usepackage{xfrac}
\usepackage[unicode,hidelinks,bookmarks=false]{hyperref}
\newcommand{\ie}{i.e.\ }
\newcommand{\cf}{cf.\ }
\newcommand{\resp}{respectively\ }
\newcommand{\Subsection}{\subsection}
\providecommand{\nobreakdash}{\nobreak\hbox{-}}
\setlength{\emergencystretch}{2em}
\multlinegap0pt
\usepackage{bm}
\usepackage{bbm}
\usepackage{dsfont}
\usepackage{xspace}
\usepackage{cancel}
\usepackage{mathtools}
\usepackage{amsthm}
\makeatletter
\@ifundefined{lemma}{\newtheorem{lemma}{Lemma}}{}
\makeatother
\newcommand{\I}{\mathbb I}
\newcommand{\R}{\mathbb{R}}
\title[Improved Power Laws for the Favard Length Problem in All Dim]{Improved Power Laws for the Favard Length Problem in All Dimensions}
\author{Caleb Marshall}
\date{Source: arXiv:2509.02882v2 (CC BY 4.0)}
\begin{document}
\maketitle

\noindent\emph{Source and licence.} Improved Power Laws for the Favard Length Problem in All Dimensions. Version arXiv:2509.02882v2 (CC BY 4.0), distributed under the Creative Commons Attribution 4.0 International licence, as recorded in the arXiv licence metadata for this exact version and shown on the abstract page https://arxiv.org/abs/2509.02882v2. Full rights notice: licenses/arxiv-2509.02882-CC-BY-4.0.txt.\par\smallskip

\noindent\textit{Editorial scope.} Counted unit: the Lemma whose source label is \texttt{lem:marked-packing} in the article (source0.tex lines 3522--3537, source label \texttt{lem:marked-packing}), delivered together with the article's own complete original proof of it (source0.tex lines 3539--3703, one \texttt{proof} environment of 165 lines). No mathematics is written, repaired or re-derived: statement and proof are the article's own text line for line. Required context retained uncounted: lma:prefixfreepacking (lines 1121--1126). Every definition, display and label of those lines is retained unchanged. Omitted from this package and not redistributed: 1--1120, 1127--3521, 3538--3538, 3704--5956. Nothing inside a retained excerpt is omitted or altered: every retained body is delivered exactly as the article's lines are written, so no omission receipt of any kind accompanies this package. Citations: the retained excerpts carry no citation command at all, so no bibliography entry of the article is transcribed and no reference is redistributed. Reference adaptation: none was needed and none was made; every reference inside the delivered unit resolves inside it, so no unresolved reference or citation is produced. Printed slips kept literal: no printed word, symbol, hypothesis or number of the counted statement or of its proof was altered. Layout and class: the article's own class is \texttt{amsart} with packages including fontenc, geometry, microtype, graphicx, subcaption, float, array, booktabs, tabularx, comment, xcolor, amsmath, amssymb, amsfonts; the wrapper is the lane's pinned \texttt{amsart} editorial wrapper at 10pt on A4, which restates the theorem-like environments the retained text needs and adds the article's own macro definitions that the retained text needs (\texttt{\textbackslash I}, \texttt{\textbackslash R}), with extra wrapper packages (bm, bbm, dsfont, xspace, cancel, mathtools). The delivered numbering is the unit's own. Assets: no retained span contains a figure environment, a graphics-inclusion command, a PDF-page inclusion, a table or a TikZ picture; no image file is copied into this package, the package declares no asset dependency and no image file is redistributed with it. Licence: version arXiv:2509.02882v2 of this article is distributed under the Creative Commons Attribution 4.0 International licence, as recorded in the arXiv licence metadata for this exact version and shown on the abstract page https://arxiv.org/abs/2509.02882v2. A full rights notice is delivered at licenses/arxiv-2509.02882-CC-BY-4.0.txt. Characters: every retained excerpt is pure ASCII, so the delivered engine, XeLaTeX with its default Latin Modern text font, reports no missing character.
\begin{lemma}[Prefix-free packing]\label{lma:prefixfreepacking}
Let $P$ be a cylinder cube and let $\mathscr F$ be a finite prefix-free family of descendants of $P$. Then
\begin{equation}\label{eq:prefixfreepacking}
\sum_{R\in\mathscr F}\ell(R)\leq\ell(P).
\end{equation}
\end{lemma}

\begin{lemma}[Maximal marked-cylinder packing]\label{lem:marked-packing}
For every $A\geq2$, there is a prefix-free family $\mathscr M_A$ of cylinder cubes, at possibly different depths, such that the following two conditions hold simultaneously.
\begin{enumerate}
\item The sum of geometric lengths along $\mathscr M_A$ is bounded by the corresponding measure of the good set at threshold $A$; which is to say,
\begin{equation}\label{eq:marked-packing}
\sum_{R\in\mathscr M_A}\ell(R)\lesssim A\mu(A).
\end{equation}
\item If $T\geq2LA$, $x\in G_T$, $f_n(x)\geq T$, then one has the following implication
\begin{equation}\label{eq:marked-packing-2}
x \in I_{\omega} \Longrightarrow Q_{\omega}\subseteq R
\text{ for some }R\in\mathscr{M}_A, \qquad \forall \omega \in \mathcal{A}^n.
\end{equation}
\end{enumerate}


\end{lemma}

\begin{proof}
We first construct the prefix-free family $\mathscr{M}_A$; we then estimate the total geometric length as in \eqref{eq:marked-packing} and establish the implication \eqref{eq:marked-packing-2}.

Set $f_0=\mathbf1_{I_{Q_\varnothing}}$.  For each $x\in G_{2A}$, define
$$
n_A (x) : = \min \big\{1 \leq n \leq N : f_{n} (x) \geq 2 A \big\},
$$
so that $n_A (x)$ denotes the minimal depth at which a stack of height $\geq 2 A$ is observed at $x$. We then define, relative to this parameter, the family of marked cylinders by setting
$$
 \mathscr{M}^{\text{marked}}_A
 :=
 \bigcup_{x\in G_{2A}}
 \left\{
 Q\in\mathscr Q_{n_A(x)}:x\in I_Q
 \right\}.
$$
Thus, a cylinder $Q$ at depth $m$ is a \textbf{marked cylinder} (or, simply, \textbf{marked}) if there exists some $x \in G_{2A}$ such that $m = n_A (x)$ and $x$ is contained in the projection $I_Q$ of $Q$. Since every cylinder has exactly $L$ children, $f_n(x)\leq Lf_{n-1}(x)$ for every $1 \leq n \leq N$ and for every $x \in \mathbb{R}$; hence,
\begin{equation}\label{eq:first-crossing}
 2A\leq f_{n_A(x)}(x)
 \leq Lf_{n_A(x)-1}(x)<2LA, \qquad \forall x \in \mathbb{R}.
\end{equation}
We record this inequality, as it is ubiquitous in later calculations.
 
To construct the family $\mathscr{M}_A$, we proceed as follows. For each $x \in G_{2A}$, consider all words $\omega \in \mathcal{A}^{n_A(x)}$ such that $x \in I_{\omega}$; which, in our notation, is the collection
$$
\mathscr{M}_{A,x}^* : = \big\{\omega \in \mathcal{A}^{n_A(x)} : x \in I_{\omega} \big\}.
$$
We then take 
$$
\widetilde{\mathscr{M}}^*_{A} : = \bigcup_{x \in G_{2A}} \mathscr{M}_{A,x}^*
$$
and remark that each word $\omega \in \widetilde{\mathscr{M}}_{A}$ need not live at some common depth. Finally, we generate a maximal prefix-free subfamily from $\widetilde{\mathscr{M}}_{A}^*$ by letting
$$
\mathscr{M}_{A}^* := \big\{\omega \in \widetilde{\mathscr{M}}_{A}^* : \omega' \not\preceq \omega, \textrm{ for any } \omega' \in \widetilde{\mathscr{M}}^*_{A} \setminus \{\omega \}\big\}.
$$
Hence, if $\tau \in \widetilde{\mathscr{M}}_{A}^* \setminus  \mathscr{M}_{A}^*$, then there necessarily exists some $\omega = \omega(\tau) \in \mathscr{M}_{A}^*$ such that $\omega \preceq \tau$ (and, in fact, this choice of prefix is unique).

We thus take
$$
\mathscr{M}_A : = \big\{Q_{\omega} : \omega \in \mathscr{M}_A^* \big\}
$$
so that each cylinder cube $Q_{\omega} \in \mathscr{M}_A$ corresponds to a unique prefix-free word $\omega \in \mathscr{M}_A^*$. Notice that, in particular, this implies that if $Q_{\omega}, Q_{\omega'} \in \mathscr{M}_A$ are distinct cylinder cubes, then necessarily $Q_{\omega} \not\subseteq Q_{\omega'}$ and $Q_{\omega'} \not\subseteq Q_{\omega}$.

We first prove \eqref{eq:marked-packing}.  Let $\mathcal D$ be the standard closed $L$-adic grid on $\mathbb R$; so that,
$$
\mathcal{D} : = \bigcup_{k \in \mathbb{Z}} \mathcal{D}_k, \textrm{ with } \mathcal{D}_k : = \big\{[m L^{-k}, (m+1) L^{-k}] : m \in \mathbb{Z} \big\}.
$$
We will use this complete $L$-adic grid $\mathcal{D}$ to organize marked intervals of different lengths.  To this end, we first define a sub-collection
\begin{equation}\label{eq:adic-candidate}
\mathcal{C}_A : = \big\{D\in\mathcal D:D\cap I_R\neq\varnothing
 \text{ for some }R\in\mathscr M_A\text{ with }\ell(R)\geq|D| \big\},
\end{equation}
before letting
\begin{equation}\label{eq:maximal-adic}
\mathcal D_{\max} : = \big\{D \in \mathcal{C}_A : \not\exists D' \in \mathcal{C}_A \textrm{ with }  D \subsetneq D' \big\}.
\end{equation}
Thus, by the construction of this family $\mathcal{D}_{\max}$, we observe that the interiors of all such $D \in \mathcal{D}_{\max}$ are pairwise disjoint; and, moreover, their union covers all possible intervals $I_R$ with $R \in \mathscr{M}_A$.

We thus fix $D\in\mathcal D_{\max}$ and write $|D|=L^{-k}$ for some $k \in \mathbb{N}$ (so that $k$, itself, is also fixed).  Maximality and the discreteness of the geometric lengths imply
\begin{equation}\label{eq:no-larger-marked}
 \ell(R)\leq|D|,
 \text{ whenever }R\in\mathscr M_A\text{ and }I_R\cap D\neq\varnothing.
\end{equation}
Indeed, this holds since, if $\ell(R)>|D|$, then the $L$-adic parent of $D$ would satisfy \eqref{eq:adic-candidate}. The same argument shows that some $R_D\in\mathscr M_A$ meeting $D$ satisfies $\ell(R_D)=|D|$. For every $R\in\mathscr M_A$ satisfying $I_R \cap D \neq \emptyset$, let $P_k(R)\in\mathscr Q_k$ denote its depth-$k$ ancestor. That such an ancestor exists is implied by \eqref{eq:no-larger-marked}. 

Suppose first, for contradiction, that there are strictly greater than $4LA$-many distinct cylinders $P_k(R)$. Then, each of their associated intervals, which have length $\sigma_\theta|D|\geq|D|$, must meet $D$; and in particular, one such interval must contain at least one of the two endpoints $y_-, y_+$ of $D$. In particular, 
\[
\max \{f_k(y_-), f_k(y_+)\} \geq\frac12\#\{P_k(R):I_{P_k(R)}\cap D\neq\varnothing\}>2LA,
\]
and we let $y$ denote one such maximal-stacking endpoint.
In particular, this shows that if
$$
\#\{P_k(R):I_{P_k(R)}\cap D\neq\varnothing\} > 4LA,
$$
then
\begin{equation}\label{eq:contradiction-index}
m :=\min \big\{1 \leq n \leq N : f_n (y) \geq 2 A\big\} \leq k - 1.
\end{equation}
Now, let $Q \in \mathscr{Q}_m$ satisfy $y \in I_Q$. Since we have established that $f_{m-1} (y) < 2 A \leq f_m (y)$, there necessarily exists some $R' \in \mathscr{M}_A$ such that $Q \subseteq R'$, and so, in particular, $\ell (R') \geq \ell (Q) > |D|$, since $m < k$. Moreover, we know that $y \in I_Q \subseteq I_{R'}$, and so $I_{R'} \cap D \neq \emptyset$. Hence, we have contradicted \eqref{eq:no-larger-marked}, and thus proven that
\begin{equation}\label{eq:minimal-depth-ancestor-bound}
\# \{P_k (R) : I_{P_k(R)} \cap D \neq \emptyset \} \leq 4LA.
\end{equation}

Now, continuing to keep the $L$-adic interval $D \in \mathcal{D}_{\max}$ fixed, we thus define a sub-family
$$
\mathscr{P}_D : = \big\{P_k (R) : R \in \mathscr{M}_A, I_R \cap D \neq \emptyset \} \subseteq \mathscr{Q}_k.
$$
For each $P\in \mathscr{P}_D$, the members $R \in \mathscr M_A$ contained in $P$ are prefix-free descendants; and, hence, the $L$-ary Kraft inequality (which we have stated in the form of Lemma \ref{lma:prefixfreepacking}) gives
$$
 \sum_{\substack{R\in\mathscr M_A\\R\subset P}}\ell(R)\leq\ell(P)=|D|.
$$
By combining this estimate with \eqref{eq:minimal-depth-ancestor-bound}, we thus obtain the inequality
\begin{equation}\label{eq:local-packing}
 \sum_{\substack{R\in\mathscr M_A\\I_R\cap D\neq\varnothing}}\ell(R)
 \leq4LA|D|.
\end{equation}

Now, choose a cylinder $R_D \in \mathscr{M}_A$ satisfying $I_{R_D} \cap D \neq \emptyset$ and $\ell (R_D) = |D|$. That such a cylinder must exist follows from the definitions of sub-families $\mathcal{C}_A$ and $\mathcal{D}_{\max}$ given in \eqref{eq:adic-candidate} and \eqref{eq:maximal-adic}, respectively. Then, there necessarily exists a point $x_D\in I_{R_D}$ at which at least $\lceil2A\rceil$ depth-$k$ projected cylinders overlap.  At least half of these projections extend from $x_D$ by $\sigma_\theta|D|/2$ on the same side.  Since
$\lceil2A\rceil/2\geq A$, there is a one-sided interval $J_D$ simultaneously satisfying
\begin{equation}\label{eq:charge-interval}
 |J_D|= \frac{|I_{R_D}|}{2}=\frac{\sigma_\theta}{2} \lvert D \rvert, \textrm{ and also }
 J_D\subset G_A.
\end{equation}
We thus let $C_D \subseteq G_A$ be the connected component containing $J_D$.  Since $I_{R_D}$ meets $D$, there exists a dimensional constant $c_d \geq 1$ such that
\begin{equation}\label{eq:component-enlargement}
 D\subset c_d C_D,
\end{equation}
where the right-hand side denotes the concentric enlargement by the factor $c_d$.  Indeed, the interval containment in \eqref{eq:component-enlargement} follows since
$$
 |C_D|\geq|J_D|=\frac{\sigma_\theta}{2}|D|\geq\frac{|D|}{2},
$$
while also one has
$$
\operatorname{dist}(D,C_D)\leq |I_{R_D}|\leq\sqrt d\,|D|.
$$


For each connected component $C$ of $G_A$, we thus define
$$
 \mathcal D(C):=\{D\in\mathcal D_{\max}:C_D=C\}.
$$
The intervals in $\mathcal D_{\max}$ have disjoint interiors, and
\eqref{eq:component-enlargement} places every member of $\mathcal D(C)$ inside
the same fixed enlargement $c_dC$.  Hence
$$
 \sum_{D\in\mathcal D(C)}|D|
 =\left|\bigcup_{D\in\mathcal D(C)}D\right|
 \leq c_d |C|.
$$
The connected components of $G_A$ are pairwise disjoint and their lengths sum to $\mu(A)$. Hence, summing over these connected components thus gives
\begin{equation}\label{eq:adic-packing}
 \sum_{D\in\mathcal D_{\max}}|D|
 \leq c_d 
 \sum_{C\in\operatorname{Comp}(G_A)}|C|
 =c_d \, \mu(A).
\end{equation}
As every $R\in\mathscr M_A$ meets at least one $L$-adic interval in $\mathcal D_{\max}$, we may thus sum over \eqref{eq:local-packing} while also summing over $D\in\mathcal D_{\max}$ and use \eqref{eq:adic-packing} to obtain
$$
 \sum_{R\in\mathscr M_A}\ell(R)
 \lesssim A\sum_{D\in\mathcal D_{\max}}|D|
 \lesssim A\mu(A),
$$
which is \eqref{eq:marked-packing}.

Finally, suppose that $T\geq 2LA$, $x\in G_T$, and
$f_n(x)\geq T$.  Set $m:=n_A(x)$.  Since
$$
2A\leq f_m(x)<2LA\leq T\leq f_n(x),
$$
we have $m<n$.  Let $Q=Q_\omega\in\mathscr Q_n$ be any
cylinder over $x$ at depth $n$ (so that $x\in I_Q$) and let
$
P \in \mathscr Q_m
$
denote its ancestor at depth $m$.  Since $Q\subseteq P$, we have
$I_Q\subseteq I_P$, and therefore $x\in I_P$.  By the definition
of $m=n_A(x)$, every cylinder at depth $m$ whose projection contains
$x$ is marked; and so, in particular, $P$ is a marked cylinder.  Every marked cylinder is contained in a member of $\mathscr M_A$, so there is
some $R\in\mathscr M_A$ such that
\[
Q\subseteq P\subseteq R.
\]
Thus every cylinder at depth $n$ over $x$ is contained in a member of
$\mathscr M_A$, proving the final assertion.
\end{proof}

\end{document}
